\section{Результаты счёта и времена получения}
\label{sec:results}

Опорные запуски выполнены при $r=m=8$, $s_h=s_w=8$, $p=16$, $\beta=1$, $T_{\max}=80$. Результаты приведены в таблице~\ref{tab:results}. Время обучения включает все эпохи до достижения порога или лимита; время сборки изображения и расчёта метрик учитывается отдельно.

\begin{table}[H]
\centering
\caption{Опорные запуски (без перекрытия, $p=16$, $\beta=1$)}
\label{tab:results}
\small
\begin{tabular}{@{}lcccccccc@{}}
\hline
Изображение & $Q$ & $K$ & $Z$ & $T$ & $E$ & $\mathrm{MSE}_n$ & PSNR, дБ & $t$, с \\
\hline
Плавные переходы & 1024 & 1{,}263 & 10{,}1 & 80 & 3{,}13 & 0{,}000027 & 51{,}64 & 1{,}71 \\
Контуры & 1024 & 1{,}263 & 10{,}1 & 80 & 812{,}14 & 0{,}010536 & 25{,}94 & 6{,}17 \\
Текстура & 1024 & 1{,}263 & 10{,}1 & 80 & 998{,}65 & 0{,}009688 & 26{,}16 & 8{,}15 \\
\hline
\end{tabular}
\end{table}


В опорных запусках использован строгий порог $\varepsilon=10^{-12}$, который за $T_{\max}=80$ эпох не был достигнут ни для одного изображения, поэтому все три обучения остановлены по лимиту эпох. Наименьшая итоговая ошибка получена для изображения с плавными переходами: такой сигнал хорошо описывается небольшим числом линейных направлений. Для контуров и текстуры остаточная ошибка выше: резкие границы и высокочастотные детали требуют большего числа компонент, чем $p=16$.

Коэффициент по методике $Z\approx10{,}1$ превышает единицу во всех запусках. Байтовый коэффициент $K\approx1{,}26$ также больше единицы, поскольку при $p=16$ объём кодов и декодера в \texttt{float64} остаётся меньше объёма исходных \texttt{uint8}-данных. При увеличении $p$ обе величины уменьшаются, причём $K$ приближённо пропорционален $1/p$; поэтому именно $K$ даёт более строгую оценку степени сжатия.

Время обучения пропорционально $Q$, $n$ и $p$. При $Q=1024$, $n=192$, $p=16$ один опорный запуск занимает от 2 до 8~с в зависимости от изображения; серии A--D, содержащие десятки запусков, выполняются за единицы минут.
